\section{Double counting between layers}\label{sec:doblecuenta}

So far we have counted repairs starting from a crossing in layer $h$. To use this count to study $f_h(n)$, we need to relate it to the
sizes $N_h$ and $N_{h-1}$. We do this by pairing each forward repair with the operation that undoes it. Fix $h\ge1$ and let $q=n-h$ and
$m=q+1=n-h+1$.

\begin{definicion}[Forward repairs and reverse operations]\label{def:reparaciones} A \emph{forward repair} starts from $P\in\Om_h$,
chooses an empty crossing $(r,c)$ and a queen $(i,j)$ in $P$, and replaces that queen with $(r,j)$ and $(i,c)$ if the resulting
placement is legal. A \emph{reverse operation} starts from $P'\in\Om_{h-1}$, chooses an ordered pair of distinct queens $(r,j),(i,c)$
in $P'$, and replaces them with $(i,j)$ if the resulting placement is legal. The resulting empty crossing is $(r,c)$. \end{definicion}

In Figure~\ref{fig:reparacion-tablero}(b), the first queen in the ordered pair is $(2,0)$ and the second is $(0,1)$. To undo the
repair, we take the row of the second queen and the column of the first, restoring the queen at $(0,0)$. The row of the first queen and
the column of the second are empty again and meet at the original crossing $(2,1)$.

The two orientations of a pair correspond to the two possible replacements, so there is no further factor of two. In a legal placement
the paired queens occupy different rows and columns and do not share a diagonal, so their eight lines are distinct. The two diagonals
of $(i,j)$ are distinct from these eight lines, and neither removed queen attacks $(i,j)$ diagonally. Thus the reverse operation is
legal exactly when the diagonal and antidiagonal of $(i,j)$ are empty in $P'$. Each legal forward repair corresponds to a unique legal
reverse operation, and conversely.

For a bounded weight $w(r,c)$ that assigns the same value to an operation and its reverse, define \[
 D_w(P)=\sum_{r\in R_\vac,\,c\in C_\vac}w(r,c)\,A^{rc}(P),\qquad
 I_w(P')=\sum_{\substack{(r,j),(i,c)\in P'\\ \text{distinct}}}w(r,c)\,V_{D(i,j)}V_{A(i,j)},
\] where $V_{D(i,j)}$ and $V_{A(i,j)}$ indicate that the diagonal and antidiagonal of $(i,j)$ are empty, respectively. For $w\equiv1$
write $d_{h,1}=D_1$ and $d_h^\inv=I_1$.

\begin{lema}[Double counting]\label{lem:doblecuenta} For every bounded weight $w$, \begin{equation}\label{eq:doblecuenta}
 N_h\,\E_{\mu_h}D_w=N_{h-1}\,\E_{\mu_{h-1}}I_w ,
\end{equation} and $|D_w|/(h^2n)\le\|w\|_\infty$, $|I_w|/n^2\le\|w\|_\infty$. In particular, with $\overline d_h=d_{h,1}/(h^2q)$ and
$\overline d{}^{\,\inv}_h=d_h^\inv/(mq)$, \[
 0\le\overline d_h,\overline d{}^{\,\inv}_h\le1,\qquad
 \E_{\mu_h}\overline d_h=f_h\,\E_{\mu_{h-1}}\overline d{}^{\,\inv}_h,\qquad
 f_h\ge\frac{\E_{\mu_h}d_{h,1}}{h^2q}.
\] If $\E_{\mu_{h-1}}I_w\ne0$, \begin{equation}\label{eq:fh_cociente}
 f_h=\frac mn\,\frac{\E_{\mu_h}[D_w/(h^2n)]}{\E_{\mu_{h-1}}[I_w/n^2]} .
\end{equation} \end{lema}

\begin{proof} Both sides of~\eqref{eq:doblecuenta} sum $w(r,c)$ over the same set of pairs (placement, legal operation), by the
bijection just described. Each crossing has at most $q$ candidate queens and there are $h^2$ crossings, so $0\le d_{h,1}\le h^2q$. A
placement in $\Om_{h-1}$ has $m$ queens and $m(m-1)=mq$ ordered pairs, so $0\le d_h^\inv\le mq$. The weighted bounds follow from $q\le
n$ and $m\le n$. Dividing~\eqref{eq:doblecuenta} by $N_h$ and using definition~\eqref{eq:fh} gives the identity between the means
and~\eqref{eq:fh_cociente}. The lower bound on $f_h$ follows from $\overline d{}^{\,\inv}_h\le1$. \end{proof}

We use the ratio~\eqref{eq:fh_cociente} only when its denominator is positive. In Section~\ref{sec:limite} positivity holds for large
$n$ because the denominator converges to a positive limit. The argument does not divide by individual degrees or require every
placement to admit a repair. Individual degrees can be zero.
